\chapter{Experiments and Results}

\section{Experimental Setup}
The implemented system was evaluated on a small set of research papers retrieved from arXiv. For each selected paper, the pipeline performed PDF download, page/section-aware parsing, chunking, embedding generation, and indexing into a local vector store. Experiments were conducted using the configured LLM and embedding model through the application interface, focusing on three functionalities: summarization, citation-grounded question answering, and side-by-side comparison of two to three papers.

\subsection{Evaluation Dataset}
A manually curated set of queries was prepared to reflect typical research tasks, including:
\begin{itemize}
    \item Factual questions answerable from specific sections (e.g., methodology, dataset, results).
    \item Summarization requests for individual papers.
    \item Comparative questions requiring synthesis across two or three papers (e.g., differences in evaluation methods or reported performance).
\end{itemize}
The evaluation scale was intentionally limited, consistent with the project's scope and the soft-evaluation nature noted in the topic definition.

\section{Evaluation Methodology}
Since large-scale ground-truth annotations were not available, evaluation combined manual inspection with qualitative LLM-as-judge style criteria \cite{b13}.

\subsection{Criteria}
\begin{itemize}
    \item \textbf{Relevance:} Whether the response addresses the user query or requested comparison aspect.
    \item \textbf{Faithfulness/Groundedness:} Whether claims are supported by retrieved context from the source papers.
    \item \textbf{Citation Quality:} Correctness and usefulness of cited section and page references for locating evidence.
    \item \textbf{Summary Quality:} Coverage of key contributions, clarity, and conciseness.
    \item \textbf{Comparison Quality:} Accuracy of aligned similarities/differences and completeness across selected papers.
\end{itemize}

\subsection{Scoring Approach}
Responses can be rated on a simple ordinal scale (e.g., 1--5 or Low/Medium/High) per criterion. Where LLM-as-judge is used, a separate evaluation prompt assesses generated outputs against the retrieved evidence and the user query. Manual verification is recommended for citation checks by inspecting the referenced section/page in the original PDF.

\section{Results and Observations}
Table~\ref{tab:eval_results} provides a template for reporting evaluation outcomes. Actual scores can be populated based on recorded runs.

\begin{table}[H]
\centering
\caption{Sample evaluation template for system outputs}
\label{tab:eval_results}
\renewcommand{\arraystretch}{1.2}
\begin{tabular}{|p{3.2cm}|p{2.2cm}|p{2.2cm}|p{2.4cm}|p{3.0cm}|}
\hline
\textbf{Task / Query Type} & \textbf{Relevance} & \textbf{Faithfulness} & \textbf{Citation Quality} & \textbf{Remarks} \\
\hline
Single-paper QA (factual) & TBD & TBD & TBD & To be filled from evaluation runs \\
\hline
Single-paper summarization & TBD & TBD & N/A or TBD & To be filled from evaluation runs \\
\hline
Two-paper comparison & TBD & TBD & TBD & To be filled from evaluation runs \\
\hline
Three-paper comparison & TBD & TBD & TBD & To be filled from evaluation runs \\
\hline
\end{tabular}
\end{table}

\subsection{Qualitative Observations}
\begin{itemize}
    \item \textbf{Citation-Grounded QA:} Metadata-aware retrieval enabled answers to include section and page references, improving traceability compared to generic responses. Citation usefulness depends on the accuracy of section detection during parsing.
    \item \textbf{Summarization:} The system produced concise summaries capturing core problem statements and methods when relevant sections were successfully retrieved. Performance varied with PDF structure and chunk quality.
    \item \textbf{Multi-Paper Comparison:} Structured prompts helped generate side-by-side contrasts across methodology and evaluation aspects for two to three papers. Comparisons were most reliable when sufficient evidence was retrieved for each paper on the same aspect.
    \item \textbf{Error Cases:} Occasional issues included incomplete extraction from complex layouts, retrieval of partially relevant chunks, and generic phrasing when evidence was sparse. These cases highlight the need for stronger parsing and verification.
\end{itemize}

\section{Result Artifacts (To Be Added)}
Screenshots of the implemented interface and sample outputs were not available during report preparation. Corresponding placeholders are provided in the Implementation chapter (Figures~\ref{fig:ui_paper_selection}--\ref{fig:ui_comparison}) and can be replaced with actual screenshots and annotated examples in future revisions.